% -----------------------------------------------------------------------------
% Set Theory lecture notes preamble
% Forest manuscript style, written with thmtools + mdframed.
% Compile with a class that has chapters, e.g. \documentclass{report} or book.
% -----------------------------------------------------------------------------

% Basic commands
\usepackage[margin=1.5in]{geometry}
\usepackage[spanish, es-noquoting]{babel}
\usepackage[T1]{fontenc}
\usepackage[utf8]{inputenc}
\usepackage{lmodern}
\usepackage{microtype}
\usepackage{amsmath, amsfonts, mathtools, amsthm, amssymb}
\usepackage[usenames,dvipsnames,svgnames]{xcolor}
\usepackage{float}
\usepackage{enumitem}
\setlist{itemsep=0.15em, topsep=0.35em}

% TikZ
\usepackage{tikz}
\usepackage{tikz-cd}
\usetikzlibrary{intersections, angles, quotes, calc, positioning}
\usetikzlibrary{arrows.meta, decorations.pathmorphing, shadows.blur}
\usepackage{pgfplots}
\pgfplotsset{compat=1.18}

% Hyperlinks
\usepackage[colorlinks=true, linkcolor=black, citecolor=black, urlcolor=black]{hyperref}

% For math environments
\usepackage{thmtools}
\usepackage[framemethod=TikZ]{mdframed}
\mdfsetup{skipabove=1em, skipbelow=0.2em}

% Palette: forest green and parchment
\definecolor{SetForest}{HTML}{183B2D}
\definecolor{SetLeaf}{HTML}{2E6B4E}
\definecolor{SetParchment}{HTML}{FBF8EF}
\definecolor{SetMint}{HTML}{F0F7F0}
\definecolor{SetGold}{HTML}{A36F16}
\definecolor{SetInk}{HTML}{1B1B18}
\definecolor{DeepGreen}{HTML}{064E3B}
\definecolor{ChalkGrey}{HTML}{F4F4F0}

\colorlet{ResultTitle}{SetForest}
\colorlet{DefinitionTitle}{SetLeaf!70!black}
\colorlet{ExampleTitle}{SetGold!80!black}
\colorlet{NotationTitle}{SetForest!85}

% Generic mdframed settings shared by the theorem styles below.
% The vertical manuscript bar is drawn with mdframed's TikZ extras.
\mdfdefinestyle{setforestbase}{%
    skipabove=0.9em,
    skipbelow=0.6em,
    innertopmargin=7pt,
    innerbottommargin=8pt,
    innerleftmargin=11pt,
    innerrightmargin=9pt,
    linewidth=0.7pt,
    roundcorner=0pt,
    splittopskip=\topskip,
    splitbottomskip=8pt,
    %nobreak=false,
    shadow=true,
    shadowsize=4pt,
    shadowcolor=black!18,
}

\mdfdefinestyle{setforestbase1}{%
    skipabove=0.9em,
    skipbelow=0.6em,
    innertopmargin=7pt,
    innerbottommargin=8pt,
    innerleftmargin=11pt,
    innerrightmargin=9pt,
    linewidth=0.7pt,
    roundcorner=0pt,
    splittopskip=\topskip,
    splitbottomskip=8pt,
    %nobreak=false,
    %shadow=true,
    %shadowsize=4pt,
    %shadowcolor=black!18,
}

\theoremstyle{definition}

\declaretheoremstyle[
    headfont=\bfseries\sffamily\color{DefinitionTitle},
    bodyfont=\normalfont,
    spaceabove=0pt,
    spacebelow=0pt,
    mdframed={nobreak,
        style=setforestbase,
        linecolor=SetLeaf,
        backgroundcolor=SetMint,
        singleextra={\draw[line width=4pt, color=SetLeaf] (O) -- (O|-P);},
        firstextra={\draw[line width=4pt, color=SetLeaf] (O) -- (O|-P);},
        middleextra={\draw[line width=4pt, color=SetLeaf] (O) -- (O|-P);},
        secondextra={\draw[line width=4pt, color=SetLeaf] (O) -- (O|-P);},
    }
]{setdefinitionbox}

\declaretheoremstyle[
    headfont=\bfseries\sffamily\color{ResultTitle},
    bodyfont=\normalfont,
    spaceabove=0pt,
    spacebelow=0pt,
    mdframed={nobreak,
        style=setforestbase,
        linecolor=SetForest,
        backgroundcolor=SetParchment,
        singleextra={\draw[line width=4pt, color=SetForest] (O) -- (O|-P);},
        firstextra={\draw[line width=4pt, color=SetForest] (O) -- (O|-P);},
        middleextra={\draw[line width=4pt, color=SetForest] (O) -- (O|-P);},
        secondextra={\draw[line width=4pt, color=SetForest] (O) -- (O|-P);},
    }
]{setresultbox}

\declaretheoremstyle[
    headfont=\bfseries\sffamily\color{ResultTitle},
    bodyfont=\normalfont,
    spaceabove=0pt,
    spacebelow=0pt,
    mdframed={
        linecolor=black!15,
        backgroundcolor=ChalkGrey,
        linewidth=0.5pt,
        %shadowcolor=black!10,
        %singleextra={\draw[line width=3pt, color=SetGold!75!black] (O) -- (O|-P);},
        %firstextra={\draw[line width=3pt, color=SetGold!75!black] (O) -- (O|-P);},
        %middleextra={\draw[line width=3pt, color=SetGold!75!black] (O) -- (O|-P);},
        %secondextra={\draw[line width=3pt, color=SetGold!75!black] (O) -- (O|-P);},
    }
]{setexamplebox}

\declaretheoremstyle[
    headfont=\bfseries\sffamily\color{NotationTitle},
    bodyfont=\normalfont,
    spaceabove=0pt,
    spacebelow=0pt,
    mdframed={
        style=setforestbase,
        linecolor=SetForest!65,
        backgroundcolor=white,
        linewidth=0.5pt,
        shadowcolor=black!10,
        singleextra={\draw[line width=3pt, color=SetForest!65] (O) -- (O|-P);},
        firstextra={\draw[line width=3pt, color=SetForest!65] (O) -- (O|-P);},
        middleextra={\draw[line width=3pt, color=SetForest!65] (O) -- (O|-P);},
        secondextra={\draw[line width=3pt, color=SetForest!65] (O) -- (O|-P);},
    }
]{setnotationbox}

\declaretheoremstyle[
    headfont=\bfseries\sffamily\color{NotationTitle},
    bodyfont=\normalfont,
    spaceabove=0pt,
    spacebelow=0pt,
    mdframed={
        %style=setforestbase,
        linecolor=SetForest!65,
        %backgroundcolor=white,
        linewidth=0.5pt,
        rightline = false, 
        leftline = false,
    }
]{setobsbox}

% Environments. Same names as your earlier notes.
\declaretheorem[style=setdefinitionbox, name=Definición, numberwithin=chapter]{definition}

\declaretheorem[style=setresultbox, name=Teorema, numberwithin=chapter]{theorem}
\declaretheorem[style=setresultbox, name=Proposición, numberwithin=chapter]{prop}
\declaretheorem[style=setresultbox, name=Lema, numberwithin=chapter]{lema}
\declaretheorem[style=setresultbox, name=Corolario, numberwithin=chapter]{colorary}

\declaretheorem[style=setexamplebox, numbered=no, name=Ejemplo]{eg}
\declaretheorem[style=setobsbox, numbered=no, name=Observación]{observation}
\declaretheorem[style=setnotationbox, numbered=no, name=Notación]{notation}
\declaretheorem[style=setexamplebox, numbered=no, name=Comentario]{remark}

% Headers and footers
\usepackage{fancyhdr}
\pagestyle{fancy}
\fancyhf{}
\fancyhead[L]{\sffamily\nouppercase{Victoria Eugenia Torroja}}
\fancyhead[R]{\sffamily\nouppercase\rightmark}
\fancyfoot[C]{\sffamily\leftmark}
\fancyfoot[R]{\sffamily\thepage}
\renewcommand{\headrulewidth}{0.4pt}
\renewcommand{\footrulewidth}{0pt}
% Generic proof name in Spanish
\renewcommand{\proofname}{Demostración}

% New commands
\newcommand{\R}{\mathbb{R}}
\newcommand{\C}{\mathbb{C}}
\newcommand{\F}{\mathbb{F}}
\newcommand{\N}{\mathbb{N}}
\newcommand{\Q}{\mathbb{Q}}
\newcommand{\Z}{\mathbb{Z}}
\newcommand{\K}{\mathbb{K}}
\newcommand{\eps}{\varepsilon}
\let\epsilon\varepsilon
\newcommand{\ds}{\displaystyle}
\newcommand{\ol}{\overline}
\newcommand{\ul}{\underline}
\newcommand{\xto}[1]{\xrightarrow{#1}}
\newcommand{\xfrom}[1]{\xleftarrow{#1}}

\DeclareMathOperator{\Ker}{Ker}
\DeclareMathOperator{\Imagen}{Im}
\DeclareMathOperator{\Hom}{Hom}
\DeclareMathOperator{\Aut}{Aut}
\DeclareMathOperator{\End}{End}
\DeclareMathOperator{\GL}{GL}
\DeclareMathOperator{\SL}{SL}
\DeclareMathOperator{\SO}{SO}
\DeclareMathOperator{\grad}{grad}
\DeclareMathOperator{\traz}{traz}
\DeclareMathOperator{\Int}{Int}
\DeclareMathOperator{\cl}{cl}
\DeclareMathOperator{\id}{id}
\DeclareMathOperator{\dist}{dist}
\DeclareMathOperator{\rank}{rank}
\DeclareMathOperator{\Span}{span}
\DeclareMathOperator{\diam}{diam}
\DeclareMathOperator{\supp}{supp}
\DeclareMathOperator{\Spec}{Spec}
\DeclareMathOperator{\Var}{Var}
\DeclareMathOperator{\Cov}{Cov}
\DeclareMathOperator{\ord}{ord}
\DeclareMathOperator{\dom}{dom}
\DeclareMathOperator{\ran}{ran}
\DeclareMathOperator{\mcd}{mcd}
\DeclareMathOperator{\mcm}{mcm}

% Set-theory helpers
\newcommand{\Pow}{\mathcal{P}}
\newcommand{\On}{\mathrm{On}}
\newcommand{\Card}{\mathrm{Card}}
\newcommand{\cf}{\operatorname{cf}}
\newcommand{\club}{\mathrm{club}}


